\section{The Ideal MPC Task}
\label{sec:ideal-mpc}

We specify the reactive task sketched in
Section~\ref{sec:overview-functionality}.  Its participants are application
pids.  A public rule chooses the successive Boolean circuits from earlier
public inputs and outputs; each circuit updates a private state, supplies
private outputs to its recipients, and produces one common public output.
Its application I/O and its separate interface to the simulator appear
on the right of Figure~\ref{fig:model-simulation}.

\subsection{Reactive tasks}
\label{sec:reactive-task}

Fix a deterministic polynomial-time rule $G$ and a public polynomial
bound $b_{\mathsf{pub}}(\lambda)$ on each public input string.
Provider $a$ in I/O round $s$ supplies a pair
\[
  x_{s,a}=(\xsec_{s,a},\xpub_{s,a})\in\mathcal X_{s,a},
  \qquad
  \mathcal X_{s,a}=\bits^{\ell_{s,a}}
      \times\bits^{\le b_{\mathsf{pub}}(\lambda)}.
\]
Write $\xsec_s=(\xsec_{s,a})_a$ and $\xpub_s=(\xpub_{s,a})_a$.
The public history selects the application circuit, which receives only
the secret input vector:
\[
  F_s=G\bigl(s,(\xpub_t,\ypub_t)_{1\le t\le s-2}\bigr),
  \qquad
  (\sigma_s,\ysec_s,\ypub_s)=F_s(\sigma_{s-1},\xsec_s).
\]
The history supplied to $G$ is empty for $s=1,2$.
Initially $\sigma_0=\epsilon$.  The circuit has $g_s$ input providers,
with secret-input lengths $\ell_{s,a}$, and $h_s$ private recipient
output bits $\ysec_s$.  The value $\ypub_s$ is a single common public
output string.  The description of $F_s$ specifies these lengths, its
public-output length, and its incoming and outgoing state lengths.
We require $G$ to produce valid circuits with compatible state lengths
along every execution.  Fix public polynomial bounds on circuit size
and on all secret-input, output, and state lengths.
Public input strings may describe later circuits directly; they need
not be passed through $F_s$.

The dependence on public inputs and outputs through round $s-2$ makes
$F_{s+1}$ available before round $s$ receives its inputs.  An application
whose next circuit depends on its most recent public inputs or output
can insert identity circuits in the even rounds:
\[
  F_{2j}(\sigma,\epsilon)=(\sigma,\epsilon,\epsilon).
\]
The odd rounds perform the application computations.  An identity round
has no application inputs or private outputs; its implementation transfers
the encrypted state to the following garbled round.

\subsection{Registration and execution}
\label{sec:mpc-execution}

Figures~\ref{fig:mpc-registration} and~\ref{fig:mpc-execution} together
define $\FMPC[G]$.  All messages carry the session identifier $\sid$.
Registration produces the same public identifier format as
Figure~\ref{fig:protocol-steps}: $\pid=(\ncek,\sigvk)$.
The ideal task uses these strings as identifiers.  It samples the NCE
registration string using $\SimNCE$ and gives the associated explanation
material to the adversary, together with the independently generated
signing key and its tape.  These are the extra registration values on
the upward arrow from $\FMPC$ to $\mathcal S$ in
Figure~\ref{fig:model-simulation}; they contain no application inputs.
In the realization proof, they let the simulator explain the keys and
tapes exposed by the implementation without choosing the application's
pids itself.

The adversary chooses when to announce each round and when to proceed.
An announcement assigns the next unused application registrations to
input providers and output recipients, in that order.  Input providers
receive their circuit and assignment at the next tick.  A subsequent
$\Proceed$ command starts a one-slot input window at the following tick.
The functionality freezes the inputs available at that tick.  During
the window, $\CommitInput$ fixes an honest provider's selected input;
$\CorruptInput$ supplies a corrupted provider's contribution.  At the
ending tick, the remaining honest contributions are fixed.  A missing
contribution has secret part $0^{\ell_{s,a}}$ and public part $\epsilon$.
Only the secret components enter the circuit; the committed public
components enter the history used by $G$.

The common public output $\ypub_s$ and corrupt recipients' private
outputs are then given to the adversary.  Honest recipients receive their
private outputs two ticks later.  This matches the label and output phases of
Figure~\ref{fig:protocol-schedule}; an intervening corruption exposes an
already-computed output before its delivery.  Once the input window is
started, these events are driven by the clock.

\begin{figure}[p]
\centering
\fbox{%
\begin{minipage}{0.95\textwidth}
\small
\textbf{The ideal task $\FMPC[G]$: registration and announcements.}

Parameters: security parameter $\lambda$, session $\sid$, the reactive
circuit rule $G$ and public-input bound $b_{\mathsf{pub}}(\lambda)$ of
Section~\ref{sec:reactive-task}, the RNCE and
signature schemes used for application registration, and a fixed RNCE
simulator $\SimNCE$.  RNCE uses the fixed delivery-message length
$\ell_{\mathsf{del}}$.

Retain the clock count $r$, initially zero; the private state
$\sigma=\epsilon$; an initially empty public history
$((\xpub_s,\ypub_s))_s$; and $\NextRound=1$.
Keep an ordered registry with each pid's corruption flag and application
record.  This record contains all its input arrivals and determined
outputs, including their times and any pending deliveries.
For each announced round retain $F_s$, the pair
$\boldsymbol{\pid}_s=(\boldsymbol{\pid}^{\mathsf{in}}_s,
\boldsymbol{\pid}^{\mathsf{out}}_s)$ of ordered input-provider and
output-recipient lists, and its announcement-delivery tick $\ReadyTick_s$.
Allocated pids remain reserved.  Execution uses Figure~\ref{fig:mpc-execution}.

\begin{description}
  \item[$\Register()$ from the adversary.]
    Sample
    \[
      (\widetilde{\ncek},\widetilde c,\tau)
          \sample\SimNCE(1^\lambda),\qquad
      (\sigvk,\sigsk)\gets\SigGen(1^\lambda;\rho_{\mathsf{sig}})
    \]
    with uniform $\rho_{\mathsf{sig}}$, and set
    $\pid=(\widetilde{\ncek},\sigvk)$.
    Reject a repeated pid.  Otherwise append it to the registry,
    with an empty application record and corruption flag false.
    Give $(\pid,\widetilde c,\tau,\sigsk,\rho_{\mathsf{sig}})$ to the
    adversary and output $\pid$ to pid $\pid$.

  \item[$\Announce(s)$ from the adversary.]
    Require that $s$ is the next unannounced round and
    $(\xpub_t,\ypub_t)_{1\le t\le s-2}$ is available.
    Compute $F_s=G\bigl(s,(\xpub_t,\ypub_t)_{1\le t\le s-2}\bigr)$.
    Require at least $g_s+h_s$ unallocated registered pids.
    Assign the first $g_s$ to the input-provider vector
    $\boldsymbol{\pid}^{\mathsf{in}}_s$ and the next $h_s$ to the
    output-recipient vector $\boldsymbol{\pid}^{\mathsf{out}}_s$,
    in registration order.  Reserve them and set $\ReadyTick_s=r+1$.
    Give $(s,F_s,\boldsymbol{\pid}_s,\ReadyTick_s)$ to the adversary.

  \item[$\Proceed(s)$ from the adversary.]
    Require $s=\NextRound$, an announced round with
    $\ReadyTick_s\leq r$, and no input-start tick yet assigned to $s$.
    Set
    \[
      \InputStart_s=r+1,\qquad
      \InputClose_s=r+2,\qquad
      \OutputTick_s=r+4.
    \]
    The clock executes these events as specified in
    Figure~\ref{fig:mpc-execution}.
\end{description}
All registered pids may be corrupted at any time.  Inputs are identified
by their sender pid; the execution and corruption handlers are in
Figure~\ref{fig:mpc-execution}.
\end{minipage}}
\caption{The ideal MPC task: registration, assignment, and starting a round.}
\label{fig:mpc-registration}
\end{figure}

\begin{figure}[p]
\centering
\fbox{%
\begin{minipage}{0.95\textwidth}
\small
\textbf{The ideal task $\FMPC[G]$: inputs, clock events, and corruption.}
Use the parameters and retained state of Figure~\ref{fig:mpc-registration}.
For each started round and provider $a$, retain the selected input pair
$\ChosenInput_{s,a}$ and committed input pair $\CommittedInput_{s,a}$,
using $\bot$ for unset values.  An eligible pair belongs to
$\mathcal X_{s,a}=\bits^{\ell_{s,a}}\times
\bits^{\le b_{\mathsf{pub}}(\lambda)}$; its default is
$(0^{\ell_{s,a}},\epsilon)$.

\begin{description}
  \item[$\Input((s,a,x))$ from a registered pid $\pid$.]
    Here $x=(\xsec,\xpub)$ is the proposed secret/public input pair.
    If this pid is corrupted, route the input to the adversary.
    Otherwise append $(r,(s,a,x))$ to its ordered input history; let $j$
    be this entry's index.  For a pair of bit strings with
    $|\xpub|\le b_{\mathsf{pub}}(\lambda)$, notify the adversary of
    $(\Input,\pid,j,s,a,\xpub)$.
    Otherwise notify it of $(\Input,\pid,j,\bot)$.
    Round, provider, and secret-length eligibility are checked at selection.

  \item[$\Input(v)$ with any other payload, from a registered pid $\pid$.]
    If this pid is corrupted, route the input to the adversary.
    Otherwise append $(r,v)$ to its ordered input history; let $j$ be
    this entry's index.  Notify the adversary of $(\Input,\pid,j,\bot)$.

  \item[$\CommitInput(s,a)$ from the adversary.]
    Require that $s$ has started, $a\in[g_s]$, $r=\InputStart_s$,
    $\pid^{\mathsf{in}}_{s,a}$ honest,
    $\ChosenInput_{s,a}\ne\bot$, and $\CommittedInput_{s,a}=\bot$.
    Set $\CommittedInput_{s,a}=\ChosenInput_{s,a}$.

  \item[$\CorruptInput(s,a,x)$ from the adversary.]
    Require that $s$ has started, $a\in[g_s]$, $r=\InputStart_s$,
    $\pid^{\mathsf{in}}_{s,a}$ corrupt,
    $x\in\mathcal X_{s,a}$, and $\CommittedInput_{s,a}=\bot$.
    Set $\CommittedInput_{s,a}=x$.

  \item[Clock tick from $\FClock$.]
    Increment $r$.  Process all events assigned to this tick, in the
    following order and in increasing $s$ within each event type.
    \begin{enumerate}
      \item If $r=\InputClose_s$, set each still-unset
        $\CommittedInput_{s,a}$ to $\ChosenInput_{s,a}$ when its provider
        is honest and this selected value exists, and to
        $(0^{\ell_{s,a}},\epsilon)$ otherwise.  Write
        $\CommittedInput_{s,a}=(\xsec_{s,a},\xpub_{s,a})$ and retain
        $\xpub_s=(\xpub_{s,a})_a$.  Compute and retain
        \[
          (\sigma',\ysec_s,\ypub_s)
            =F_s\bigl(\sigma,(\xsec_{s,a})_a\bigr),
          \qquad \sigma\gets\sigma'.
        \]
        Append $(\xpub_s,\ypub_s)$ to the public history.
        Record $(s,j,\ysec_{s,j},\OutputTick_s)$ for each recipient
        $\pid^{\mathsf{out}}_{s,j}$.  Give $(s,\xpub_s,\ypub_s)$ and
        the addressed private output bits of corrupt recipients to the
        adversary.
      \item If $r=\OutputTick_s$, output $\ysec_{s,j}$ to each still-honest
        pid $\pid^{\mathsf{out}}_{s,j}$, in increasing $j$, and mark the
        deliveries in their records.  Set $\NextRound=s+1$.
      \item For each $s$ with $r=\ReadyTick_s$, in increasing $s$, output
        \[
          (s,F_s,\boldsymbol{\pid}_s)
        \]
        to each honest $\pid^{\mathsf{in}}_{s,a}$, in increasing $a$.
        Record these outputs in the corresponding application records.
      \item If $r=\InputStart_s$, freeze each provider's input list.
        Let $j$ be the first entry whose value has the form
        $(s,a,x)$ with $x\in\mathcal X_{s,a}$.
        Set $\ChosenInput_{s,a}=x$, or $\bot$ if no such entry exists,
        and set $\CommittedInput_{s,a}=\bot$.
        Give $(\SelectedInput,s,a,j)$ to the adversary, with $j=\bot$
        when there is no selected entry.
    \end{enumerate}

  \item[On corruption of a registered pid $\pid$.]
    Mark it corrupted.  Give its application record to the adversary,
    including all input arrivals, its selected and committed input if
    defined, and all already-determined outputs with their delivery times.
\end{description}
Reject malformed commands and commands whose conditions do not hold.
\end{minipage}}
\caption{The ideal MPC task: input commitment, evaluation, and timed outputs.}
\label{fig:mpc-execution}
\end{figure}
